% Cuerpo científico recuperado y distribuido en su residencia terminal.
% Cuerpo editor-ready sin preámbulo, portada ni metadatos publicables.
% Complementa, sin sustituirlos, los cierres operatoriales y espectrales.

\section{Clôtures négatives et calibrées des équivalences apparentes}
\label{sec:hmt-final-clausuras-negativas-calibradas}

Les résultats de cette section s'appliquent après la génération conjointe
du continu. Leur antécédent causal est l'état unique produit par
APP--TRIT--TPK, l'émergence corrélative de ses cinq constructions
consubstantielles, les carrés de naturalité et le passage conjoint au terminal
et à la limite. Aucune des équivalences examinées ci-dessous ne génère ni ne
sélectionne rétrospectivement ces constructions.

\subsection{Obstructions à la canonicité à la frontière
\texorpdfstring{\(108{:}144=90{:}120\)}{108:144=90:120}}

\begin{theorem}[Il n'existe pas d'inclusion canonique
\texorpdfstring{\(6\hookrightarrow8\)}{6 à 8} sans jauge]
\label{thm:hmt-final-no-inclusion-canonica-seis-ocho}
Soient \(A\) et \(B\) des ensembles nus avec

\[
 |A|=6,
 \qquad |B|=8.
\]
Il n'existe pas de règle, naturelle par rapport à tous les réétiquetages de \(A\)
et de \(B\), qui sélectionne une injection \(A\hookrightarrow B\).
\end{theorem}

\begin{proof}
Si une injection \(\iota:A\hookrightarrow B\) était sélectionnée sans donnée
supplémentaire, la naturalité par rapport à tout \(\beta\in\operatorname{Sym}(B)\)
imposerait \(\beta\iota=\iota\). Une fois choisi \(a\in A\), il existe une permutation
\(\beta\) qui déplace \(\iota(a)\), contradiction. Le transducteur
route par route APP--Witt n'est donc pas canonique à partir des cardinaux. Une fois
fixées une octade, une inclusion et la jauge d'incidence, on obtient bien un
transducteur calibré ; le choix fait partie de son domaine.
\end{proof}

\begin{theorem}[Une fibre nue de cardinal \(28\) ne détermine pas
\(J(8,2)\)]
\label{thm:hmt-final-no-johnson-desde-cardinal}
Soit \(F\) un ensemble sans structure avec \(|F|=28\). Il n'existe pas de relation
d'adjacence non triviale sur \(F\), naturelle par rapport à toutes ses
permutations, qui produise le graphe de Johnson \(J(8,2)\).
\end{theorem}

\begin{proof}
Le groupe \(\operatorname{Sym}(F)\) est transitif sur les paires non ordonnées
d'éléments distincts. Une relation simple et invariante sous tout ce groupe
contient donc toutes les paires ou aucune. Les deux graphes résultants sont
le graphe complet et le graphe vide. En revanche, \(J(8,2)\) a \(28\) sommets et est de
degré \(12\) ; ce n'est aucun des deux. Par conséquent, la projection
calibrée \(1008\to36\) à fibres uniformes de cardinal \(28\) n'implique pas
à elle seule une octade de Witt. Une fois fixées une octade \(O\) et une bijection
\(F\simeq\binom{O}{2}\), l'adjacence de Johnson peut être transportée sur \(F\),
mais elle dépend précisément de cette jauge.
\end{proof}

Ces deux théorèmes clôturent négativement les anciennes formulations
canoniques. Ils n'affectent pas l'invariant exact
\[
 \frac{108}{144}=\frac{90}{120}=\frac68
\]
ni le défaut normalisé \(-1/12\), qui appartiennent à l'extension de rang
et non au choix d'une incidence de Witt.

\subsection{Le tour nonadique complet et le lecteur exceptionnel}

Fixons la jauge
\[
 \chi_{\mathrm{exc}}
 =(
 \prec_{\mathrm{proj}},
 \kappa_{\mathrm{Paley}},
 \mathcal D_{12},
 \iota_{\mathrm{inc}},
 \mathcal K_{\Lambda})
\]
et le lecteur calibré
\(E_{\chi_{\mathrm{exc}}}:\mathscr S_\infty
\to\mathcal I^{\mathrm{exc}}_{\chi_{\mathrm{exc}}}\).
Soit
\[
 \mathscr S_\infty\xleftarrow{s_\infty}
 \mathsf Z^{\Gamma}_{9,\infty}
 \xrightarrow{t_\infty}\mathscr S_\infty
\]
l'arête attestée d'un tour complet. Notons
\(\operatorname{ed}_0\) l'arête centrale, \(\operatorname{ev}_0\)
l'évaluation à l'état initial et \(F_9\) le décalage d'un tour.

\begin{theorem}[Factorisation exceptionnelle par l'arête attestée]
\label{thm:hmt-final-factorizacion-gamma-nueve}
Définissons
\[
 \overline E^{s}_{\chi_{\mathrm{exc}}}
 :=E_{\chi_{\mathrm{exc}}}\circ s_\infty,
 \qquad
 \overline E^{t}_{\chi_{\mathrm{exc}}}
 :=E_{\chi_{\mathrm{exc}}}\circ t_\infty.
\]
Alors
\begin{align}
 E_{\chi_{\mathrm{exc}}}\circ\operatorname{ev}_0
 &=\overline E^{s}_{\chi_{\mathrm{exc}}}
   \circ\operatorname{ed}_0,
 \label{eq:hmt-final-factorizacion-gamma-fuente}\\
 E_{\chi_{\mathrm{exc}}}\circ\operatorname{ev}_0\circ F_9
 &=\overline E^{t}_{\chi_{\mathrm{exc}}}
   \circ\operatorname{ed}_0.
 \label{eq:hmt-final-factorizacion-gamma-destino}
\end{align}
\end{theorem}

\begin{proof}
Les identités d'extrémités sont
\(s_\infty\operatorname{ed}_0=\operatorname{ev}_0\) et
\(t_\infty\operatorname{ed}_0=\operatorname{ev}_0F_9\). La composition avec
le lecteur calibré donne les deux équations.
\end{proof}

\begin{corollary}[Non-équivalence entre le secteur temporel \(+1\) et \(1\)
tour nonadique complet]
\label{cor:hmt-final-no-equivalencia-retorno-gamma}
La factorisation précédente ne peut pas être remplacée, de manière naturelle, par une
factorisation à travers le seul échantillon TRIT \(\tau=+1\).
\end{corollary}

\begin{proof}
La valeur \(\tau=+1\) enregistre le retour, la mémoire et le passé, mais oublie les données
de phase, de retenue, de route et d'extrémité qui distinguent les arêtes complètes. Deux arêtes
ayant le même échantillon temporel peuvent avoir des destinations différentes. Toute application qui
se factorise par cet échantillon identifie les deux arêtes, tandis que
\(\operatorname{ed}_0\) et ses applications source et destination les séparent. Par conséquent,
le secteur \(+1\) n'est pas équivalent à l'arête complète \(K\mapsto K+9\).
En particulier,
\[
 g(K+9)=g(K)
 \quad\text{et}\quad
 B_{K+9}\ne B_K\quad\text{en général}
\]
restent des affirmations distinctes.
\end{proof}

\begin{proposition}[Séparation calibrée face à la publication brute]
\label{prop:hmt-final-separacion-enriquecida-no-cruda}
Supposons que le lecteur exceptionnel enrichi conserve, pour chaque
générateur, le repère \(f_j\) et l'image
\(y_j=\Gamma_{f_jA_Wf_j^{-1}}(w_j)\). Alors il retrouve
\[
 w_j=\Gamma_{f_jA_Wf_j^{-1}}^{-1}(y_j)
\]
et est injectif sur le quotient par les changements de coordonnées inclus dans la
jauge. Si le lecteur brut est obtenu en oubliant les repères ou la mémoire,
\(E^{\mathrm{cr}}=\pi E^{\sharp}\), il n'est pas équivalent au lecteur enrichi
à moins que \(\pi\) ne soit injective sur \(E^{\sharp}(X)\).
\end{proposition}

\begin{proof}
La première affirmation est l'inversion explicite d'un automorphisme dans chaque
fibre. Pour la seconde, si \(\pi\) identifie deux sorties enrichies
distinctes, le lecteur brut identifie leurs antécédents et ne peut pas être
injectif. Réciproquement, l'injectivité de \(\pi\) sur l'image permet de
retrouver la sortie enrichie. Ainsi se clôt la séparation dans le domaine
calibré et se trouve rejetée l'équivalence tacite avec le codomaine brut.
\end{proof}

L'ancien carré écrit avec des symboles non typés
\(\mathrm{pub}_{\mathrm{inc}},\mathrm{coord},\mathcal O\) ne constitue pas une
équivalence mathématique : il ne définit même pas ses morphismes. Son substitut est le
cône naturel des publications archimédienne, exceptionnelle calibrée et
dimensionnelle, dont les carrés de troncature ont bien un domaine, un codomaine et une
action déclarés.

\subsection{Corridor exceptionnel et boucles du transport}

La chaîne Paley--Witt--Golay--Leech--\(V^{\natural}\)--Monstre--Moonshine
est une publication ultérieure de l'état déjà construit. Une étoile
d'incidence et une boucle du groupoïde TPK ne sont pas des objets équivalents par leur
nom : la première est un sous-objet combinatoire ; la seconde est une flèche dont
la source égale la destination, avec composition et inverse. Sans base, orientation et foncteur
de transport, il n'existe pas de comparaison typée. L'identification
automatique est donc close comme non-équivalence. Un foncteur calibré de
transport peut produire des boucles à partir d'étoiles, mais constitue un
résultat supplémentaire et non une prémisse du corridor déjà démontré.

\subsection{Canal infiniment divisible induit par le transfert TPK}

\begin{theorem}[Poissonisation UCP et racines cohérentes]
\label{thm:hmt-final-poissonizacion-ucp}
Soit \(K\) un canal UCP compatible avec les troncatures TPK\@. Pour \(t\ge0\),
définissons
\[
 P_t:=\exp\bigl(t(K-I)\bigr)
 =e^{-t}\sum_{n=0}^{\infty}\frac{t^n}{n!}K^n.
\]
Alors \((P_t)_{t\ge0}\) est un semi-groupe UCP compatible avec toutes les
troncatures et
\[
 P_1=(P_{1/m})^m
 \qquad(m\ge1).
\]
À chaque niveau fini, une dilatation minimale de Stinespring de \(P_{1/m}\)
produit une racine tensorielle cohérente ; la limite projective conserve les
marginales TPK.
\end{theorem}

\begin{proof}
Chaque \(P_t\) est une combinaison convexe convergente en norme des canaux
\(K^n\), et est donc UCP\@. Comme tous les termes sont des puissances du même
opérateur, l'identité exponentielle donne \(P_tP_s=P_{t+s}\). La compatibilité
avec les troncatures se vérifie terme à terme. En prenant \(t=1/m\), on
obtient la racine \(m\)-ième. Stinespring transforme chaque composition finie en
une contraction tensorielle ; la naturalité des canaux rend leurs
marginales cohérentes.
\end{proof}

Le théorème clôt l'existence d'un canal infiniment divisible à
mémoire TPK\@. Il n'affirme pas qu'un canal historique arbitraire coïncide avec cette
poissonisation : cette égalité exigerait de comparer leurs générateurs et reste
rejetée tant qu'une telle comparaison n'est pas fournie.

\section{Clôtures spectrales et frontières de reconnaissance externe}
\label{sec:hmt-final-clausuras-espectrales-externas}

\begin{proposition}[Équivalence calibrée, et non égalité ambiante]
\label{prop:hmt-final-equivalencia-calibrada-vnueve}
Le transport spectral du double cercle détermine une classe unitaire dans le
sous-espace cyclique engendré par son vecteur de référence. Deux représentants
ambiants \(V_9\) ne sont littéralement égaux qu'après avoir fixé une
identification unitaire du complément et une jauge compatible. Sans ces
données, l'égalité ambiante ne reste pas indéterminée : c'est une affirmation plus
forte et non équivalente à la classe spécifique du sous-espace déjà construite.
\end{proposition}

La convergence en norme des résolvantes comprimées, la récupération des
hauteurs, le critère de Riesz pour les multiplicités et la trace diagonale
RH--Moonshine sont clos dans leurs théorèmes respectifs. La trace diagonale
couple l'entier de mémoire au degré de \(V^{\natural}\) ; elle n'affirme pas que la
multiplicité de chaque zéro est la dimension d'une représentation du
Monstre. Cette identification supplémentaire n'est pas équivalente à l'existence de
la trace couplée.

Restent hors du domaine interne, sans rouvrir aucun théorème HMT :
\begin{enumerate}
 \item l'identification de la dichotomie des publications HMT analytiques à
 tous les sous-ensembles arbitraires de \(\mathbb R\) de l'univers ambiant ;
 \item la réalisation macroscopique ou thermodynamique du cristal temporel
 fini HMT déjà construit, y compris la conservation de la cohérence à grand
 volume et sa confrontation expérimentale ;
 \item la consolidation décimale des fenêtres spectrales par une
 exécution concrète d'arithmétique des intervalles ;
 \item une identification représentationnelle supplémentaire entre les multiplicités
 spectrales individuelles et les modules irréductibles du Monstre.
\end{enumerate}
Ces quatre frontières sont des reconnaissances ou des comparaisons externes, et non
 des lacunes de la construction du continu, de la clôture cardinale, de RH ni de
l'entrelacement du double cercle.

La résidence terminale conserve, sans permutation,
\[
\begin{aligned}
 \mathrm{RH}
 &\longrightarrow \text{double cercle et champ spectral}
 \longrightarrow \mathrm{Navier\text{--}Stokes}\\
 &\longrightarrow \mathrm{Yang\text{--}Mills}
 \longrightarrow \mathrm{Hodge}
 \longrightarrow \mathrm{BSD}
 \longrightarrow \text{réfutateurs}.
\end{aligned}
\]
